\section{Challenge Period Trade-offs}


The challenge period duration involves complex trade-offs between security, user experience, and capital efficiency.

\subsection{Safety Considerations}

Longer challenge periods increase safety by allowing more time for fraud detection:

\begin{equation}
P(\text{fraud detected}) = 1 - e^{-\lambda \tau}
\end{equation}

This exponential relationship suggests diminishing returns for very long periods.

\subsection{User Experience Impact}

Finality delay affects user experience:

\begin{equation}
\text{UX Score} = \frac{1}{1 + \alpha \tau}
\end{equation}

where $\alpha$ represents user sensitivity to delays.

\subsection{Capital Efficiency}

Longer periods lock capital, creating opportunity costs:

\begin{equation}
\text{Cost}_{\text{capital}} = B \times r \times \tau
\end{equation}

where $B$ is the bond amount and $r$ is the risk-free rate.

\subsection{Adaptive Period Approach}

Lux implements adaptive challenge periods based on stake and history:

\begin{equation}
\tau_{\text{adaptive}} = \tau_{\text{base}} \times f(\text{stake}, \text{history}, \text{value})
\end{equation}

where the function $f$ considers:
\begin{itemize}
    \item Stake ratio: Higher stakes enable shorter periods
    \item Historical behavior: Good actors earn shorter periods
    \item Transaction value: Higher values require longer periods
\end{itemize}

